Our experimental design tests four predictions via a phased protocol: (P1) correlation existence beyond random chance, (P2) stable failure mode clusters, (P3) scale-invariant patterns, and (P4) intervention targeting. We report results for P1-P2; P3-P4 were blocked or deferred due to clustering quality failure (Section 5).

\subsection{Research Questions}

\textbf{RQ1 (P1):} Do trustworthiness failures correlate across dimensions beyond random chance?
\begin{itemize}
\item \textbf{Hypothesis:} Spearman $r > 0.3$ and $p < 0.01$ after Bonferroni correction for $\geq 70\%$ of benchmark pairs (proof-of-concept: $\geq 1$ pair).
\item \textbf{Rationale:} Medium effect size $r > 0.3$ indicates shared root causes; stronger correlations suggest tighter coupling or dimensional redundancy.
\end{itemize}

\textbf{RQ2 (P2):} Do correlated failures cluster into distinct, stable modes?
\begin{itemize}
\item \textbf{Hypothesis:} Hierarchical clustering yields silhouette score $> 0.5$ AND bootstrap consistency $\geq 80\%$.
\item \textbf{Rationale:} Well-separated (silhouette $> 0.5$) and stable (bootstrap $\geq 80\%$) clusters would validate interpretable failure mode taxonomy.
\end{itemize}

\textbf{RQ3 (P3):} Are correlation patterns scale-invariant across model sizes?
\begin{itemize}
\item \textbf{Hypothesis:} Mantel test $r > 0.7$ comparing correlation matrices across size strata.
\item \textbf{Rationale:} Scale invariance indicates fundamental property rather than architecture/size artifact.
\end{itemize}

\textbf{RQ4 (P4):} Do interventions targeting one cluster improve all cluster benchmarks without degrading others?
\begin{itemize}
\item \textbf{Hypothesis:} Cohen's $d > 0.5$ for cluster benchmarks, $d < 0.1$ for non-cluster benchmarks.
\item \textbf{Rationale:} Cluster-specific intervention effects validate actionable taxonomy.
\end{itemize}

\subsection{Dataset Construction}

\textbf{Model Corpus ($n=20$):} Public benchmark results for LLMs spanning three size strata:
\begin{itemize}
\item \textbf{Small ($<$1B params, $n=6$):} GPT-2 variants, Phi-2, TinyLlama, smaller open-source models
\item \textbf{Medium (1-10B, $n=9$):} LLaMA-7B, Mistral-7B, GPT-3.5-turbo, Claude-instant, mid-size variants
\item \textbf{Large ($>$10B, $n=5$):} GPT-4, Claude-3-Opus, LLaMA-70B, large proprietary models
\end{itemize}

Size distribution intentionally oversamples medium models (45\% of corpus) where public data availability is highest. Stratification enables confound control: if correlations persist within size-homogeneous groups, coupling cannot be attributed solely to size-related performance trends.

\textbf{Benchmark Scores:} Aggregated from published leaderboards, model cards, and research papers:
\begin{itemize}
\item TrustfulQA: MC1 accuracy (\% correct on single-answer questions)
\item AdvBench: Defense rate (\% of adversarial attacks successfully blocked)
\item BOLD: Composite fairness score (1 - max demographic disparity across sentiment/regard metrics)
\end{itemize}

All scores normalized to [0,1] via min-max scaling. Missing data handling: Models without all three benchmark scores excluded (reduced candidate pool from 35 to 20 models).

\subsection{Evaluation Metrics}

\textbf{Primary Metrics (RQ1 - Correlation):}
\begin{itemize}
\item \textbf{Spearman's $\rho$:} Rank correlation coefficient (range -1 to 1). Non-parametric, robust to non-linear monotonic relationships.
\item \textbf{p-value:} Permutation test significance with Bonferroni correction ($\alpha = 0.01 / 3 = 0.0033$ for 3 pairwise tests).
\item \textbf{Effect size interpretation:} $r > 0.1$ small, $r > 0.3$ medium, $r > 0.5$ large (Cohen's conventions).
\end{itemize}

\textbf{Secondary Metrics (RQ2 - Clustering):}
\begin{itemize}
\item \textbf{Silhouette score:} Separation quality (range -1 to 1). Values $> 0.5$ = well-separated, 0.3-0.5 = acceptable, $< 0.3$ = poor.
\item \textbf{Bootstrap consistency:} \% of 1000 bootstrap iterations where benchmark pairs cluster together. $\geq 80\%$ indicates stable clusters.
\item \textbf{Cophenetic correlation:} Dendrogram distance preservation quality. $> 0.7$ acceptable.
\end{itemize}

\textbf{Tertiary Metrics (RQ3 - Scale Invariance, not evaluated):}
\begin{itemize}
\item \textbf{Mantel test $r$:} Matrix correlation comparing correlation structure across strata. $> 0.7$ = scale-invariant.
\end{itemize}

\textbf{Quaternary Metrics (RQ4 - Interventions, not evaluated):}
\begin{itemize}
\item \textbf{Cohen's $d$:} Pre/post intervention effect size. $d > 0.5$ = medium effect, $d > 0.8$ = large.
\end{itemize}

\subsection{Baseline Comparisons}

\textbf{Null Model (Permutation Test):} Generate null distribution by shuffling benchmark scores 1000 times, recomputing correlations for each permutation. Observed correlations compared to 99th percentile of null distribution ($\alpha = 0.01$ threshold).

\textbf{Comparison to Existing Frameworks:}
\begin{itemize}
\item \textbf{HELM aggregation:} Reports separate scores without correlation analysis. Our approach extends this by discovering latent structure.
\item \textbf{Single-dimension evaluation:} TrustfulQA, AdvBench, BOLD analyzed in isolation. We test whether scores correlate across these benchmarks.
\end{itemize}

\subsection{Experimental Procedure}

\textbf{Phase 1 (H-E1): Correlation Existence}
\begin{enumerate}
\item Collect 20 models $\times$ 3 benchmark scores (60 observations).
\item Compute Spearman $\rho$ for 3 benchmark pairs: TrustfulQA $\leftrightarrow$ AdvBench, TrustfulQA $\leftrightarrow$ BOLD, AdvBench $\leftrightarrow$ BOLD.
\item Run permutation test (1000 iterations) for each pair.
\item Apply Bonferroni correction: reject null if $p < 0.0033$.
\item Count significant pairs (success: $\geq 1$ for PoC, $\geq 2$ for full validation).
\item Stratified analysis: Repeat within small/medium/large subgroups.
\end{enumerate}

\textbf{Phase 2 (H-M1): Cluster Stability}
\begin{enumerate}
\item Convert correlation matrix to distance matrix: $d = 1 - r$.
\item Apply Ward linkage hierarchical clustering.
\item Evaluate $k=2$ clusters ($k \geq 3$ not computable with $n=3$ benchmarks).
\item Compute silhouette score for $k=2$ solution.
\item Bootstrap resampling: 1000 iterations, track cluster co-occurrence.
\item Gate decision: Proceed to Phase 3 only if silhouette $> 0.5$ AND bootstrap $\geq 80\%$.
\end{enumerate}

\textbf{Phase 3 (H-M4): Scale Invariance (blocked)}
\begin{itemize}
\item Planned: Stratified clustering + Mantel test on strata correlation matrices.
\item Status: Not executed due to Phase 2 gate failure (MUST\_WORK).
\end{itemize}

\textbf{Phase 4 (Intervention): Validation (deferred)}
\begin{itemize}
\item Planned: Apply calibration training to highest-silhouette cluster, measure pre/post Cohen's $d$.
\item Status: Deferred due to absence of validated clusters (depends on Phase 2).
\end{itemize}

\subsection{Implementation Details}

\textbf{Software:} Python 3.10, scipy.stats for correlations, sklearn.cluster for hierarchical clustering, sklearn.metrics for silhouette scores.

\textbf{Reproducibility:} All data, code, and intermediate results available at [repository link]. Random seed fixed for permutation tests and bootstrap resampling (seed=42).

\textbf{Compute:} Analysis completed on single CPU (correlation + clustering: $\sim$5 minutes total runtime). No GPU required.

\subsection{Why This Design Tests the Hypothesis}

The four-phase protocol directly maps to hypothesis predictions:
\begin{itemize}
\item \textbf{P1 validated} $\rightarrow$ Failures correlate (shared root causes exist)
\item \textbf{P2 validated} $\rightarrow$ Correlations organize into distinct modes (actionable taxonomy)
\item \textbf{P3 validated} $\rightarrow$ Patterns generalize across scales (not architecture artifact)
\item \textbf{P4 validated} $\rightarrow$ Taxonomy enables targeted interventions (practical utility)
\end{itemize}

Conversely, failures at each phase falsify specific claims:
\begin{itemize}
\item \textbf{P1 fails} $\rightarrow$ Independent dimensions (abandon hypothesis)
\item \textbf{P2 fails} $\rightarrow$ No distinct taxonomy despite correlations (explore alternative clustering methods or expand benchmarks)
\item \textbf{P3 fails} $\rightarrow$ Scale-specific patterns (limit generalization scope)
\item \textbf{P4 fails} $\rightarrow$ Taxonomy not actionable (theoretical contribution only)
\end{itemize}

Actual results: P1 far exceeded expectations ($r > 0.99$ instead of $r > 0.3$), P2 failed gate (silhouette = 0.274 despite 100\% bootstrap), P3/P4 untested. This outcome validates tight coupling while refuting distinct failure mode taxonomy, revealing 3-benchmark design limitation rather than falsifying correlation-based approach.
